% ---------------------------------------------------------------------------
% Formatação do currículo. Portado do template LaTeX antigo (Overleaf).
% Nada aqui precisa ser tocado no dia a dia: para editar o conteúdo,
% mexa apenas nos arquivos cv-en.qmd e cv-pt.qmd.
% ---------------------------------------------------------------------------

\usepackage{fontawesome5}
\usepackage{titlesec}
\usepackage{enumitem}
\usepackage{xcolor}
\usepackage{array}
\usepackage{needspace}
\usepackage{etoolbox}

\setlength{\parindent}{0pt}
\setlength{\parskip}{0.15cm}
\pagestyle{empty}

% --- Nome + linha de contato ------------------------------------------------

% Versaletes sem negrito: é o que o template antigo produzia na prática,
% já que a família Latin Modern não tem versaletes em negrito.
\newcommand{\cvname}[1]{%
  \begingroup
  \centering
  \huge\scshape #1\par
  \vspace{0.4ex}
  \rule{\textwidth}{1.2pt}\par
  \endgroup
  \vspace{0.3cm}%
}

% O título do YAML (title:) vira o cabeçalho do currículo.
\makeatletter
\renewcommand{\maketitle}{\cvname{\@title}}
\makeatother

% Linha de contato: o bloco ::: contact do .qmd cai aqui.
\newenvironment{cvcontact}
  {\par\centering\small}
  {\par\vspace{0.3cm}}

% Separador entre os itens da linha de contato.
\newcommand{\cvsep}{\hspace{0.85em}}

% --- Seções (\section) ------------------------------------------------------

\titleformat{\section}
  [hang]
  {\Large\bfseries}
  {}{0pt}
  {}
  [\vspace{-0.25cm}\rule{\textwidth}{0.8pt}]
\titlespacing*{\section}{0pt}{0.45cm}{0.2cm}

% --- Subseções (\subsection): agrupamentos dentro de uma seção --------------

\titleformat{\subsection}
  [hang]
  {\large\bfseries}
  {}{0pt}
  {}
\titlespacing*{\subsection}{0pt}{0.35cm}{0.12cm}

% --- Entradas (\subsubsection): cada emprego, curso, publicação etc. --------

\titleformat{\subsubsection}
  [hang]
  {\normalsize\bfseries}
  {\faAngleRight}{0.45em}
  {}
\titlespacing*{\subsubsection}{0pt}{0.32cm}{0.06cm}

% Evita que uma entrada fique órfã (ou quase) no pé da página.
\pretocmd{\subsubsection}{\Needspace*{4\baselineskip}}{}{}
\widowpenalty=10000
\clubpenalty=10000

% --- Etiqueta de data (caixa preta à direita) -------------------------------
% Gerada pelo filtro _cv.lua a partir de [texto]{.date} num título de entrada.

\newcommand{\cvdate}[1]{%
  \hfill\colorbox{black}{\hspace{0.1cm}\textcolor{white}{\small #1}\hspace{0.1cm}}%
}

% --- Listas -----------------------------------------------------------------

\setlist[itemize]{leftmargin=1.1em, topsep=0.08cm, itemsep=0.04cm, parsep=0pt}
\setlist[itemize,1]{label=\textendash}

% --- Tabelas (referências) --------------------------------------------------

\usepackage{booktabs}
\renewcommand{\arraystretch}{1.15}
\AtBeginEnvironment{longtable}{\small}
